\section{Limitations and Future Directions}
The corpora are small. The two Chinese audio windows appear to contain nearly the same script, and their recording origin is undocumented. The English holdout and its safe times are authored test contracts, not independent human judgments of comedic appropriateness. Allowed alternatives improve specificity of evaluation but do not make the labels a perceptual ground truth. Negative controls are necessary checks, not evidence of reliable abstention across ordinary speech in general.

The revised model and policy are selected using a small development set. The fresh holdout prevents direct tuning on its outputs, but a single authored holdout cannot establish generalization across performers, accents, languages, or spontaneous humor. In particular, lexical ASR stability cannot establish that a later clause will not change a phrase's meaning. Prosody, deliberate pauses, audience laughter, and overlapping speech are not represented adequately by constant word clocks.

Scheduling replay assumes one nonpreemptive worker and fixed per-request demands. Scaled overload scenarios are not physical GPU contention measurements. Short controlled payloads isolate queue behavior but do not evaluate semantic selection. All-skip settings remain failures even when stale displays disappear. A deployed pipeline should be evaluated with simultaneous ASR and generation, measured transport and texture costs, and workloads extending across utterance boundaries.

The next implementation step is a bounded dynamic-text renderer or approved-asset selector with explicit candidate, commit, cancel, and presentation acknowledgments. Audio-clock alignment should include offset uncertainty. A synchronized microphone and optical display capture can independently identify the physical word-end-to-first-frame boundary. The present show log and software audio clock cannot provide that measurement.

The original audience questions also remain open. Actual authorized sessions are required to evaluate funniness, presence, distraction, readability, and comfort. A subsequent study should compare timing policies while holding content and visual effects constant, and ask annotators to identify safe reveal windows and acceptable alternative emphasis. Engineering improvements in the present pilot should guide those experiments, not substitute for them.

\section{Conclusion}
VRStandup now has an experimentally inspectable backend for proposing expressive text from partial transcripts, validating cached candidates, and controlling task freshness. The component study establishes reliable execution on the available RTX 4090 and documents a substantial improvement in output compliance after model selection. The paired streaming experiment evaluates the revised commitment policy on a fresh authored holdout, including negative and revision controls. The controlled overload test demonstrates that latest-only scheduling improves fresh-supported coverage and that deadline guards reduce additional wasted work without losing that fresh output count in the primary setting. These are concrete prototype and systems-evaluation results with explicit boundaries; connecting the backend to a headset and establishing audience benefit remain future work.
